Kubernetes is the foundation for providing the unified interface and will provide the standardized \gls{api} for the software implemented.
Therefore, the following section will provide an introduction to the core concepts of Kubernetes, using promise theory as a way to explain functionalities.
First, there will be an introduction to Kubernetes itself, and how it operates, followed by the concept of controllers.
Following that will be an introduction to operators and their concepts, including how they utilize Kubernetes objects.

\subsection{Core Kubernetes concepts}
Kubernetes works by storing objects in a key-value storage called etcd, which is accessed through the Kubernetes API server. The cluster then uses the data stored to determine a desired state of the cluster, and will try to move the actual state of the cluster closer to the desired state.
This maps closely to the concept of a promise being made.
The cluster promises to provide the desired state.
Whether it is successful in doing so is something it reports back to the user consuming the promise.
\parencite{poulton2025kubernetes, k8sclustercomponents}

\subsection{Controllers and promises in Kubernetes}
Controllers manage a specific type of object in Kubernetes, and may watch any number of other objects.
They are agents making promises about what they can do, and other agents consume them.
\parencite{k8sclustercomponents, promise-theory, poulton2025kubernetes}
A controller uses a reconcile loop in order to move the actual or observed state of the cluster closer to the desired state as defined by a manifest.
Assuming a \texttt{Deployment} object specifies that three instances of a \texttt{Pod} with an nginx container should be running, the controller responsible for \texttt{Deployment} objects will reconcile the object and assess its condition to ensure that there are actually three instances running.
\parencite{poulton2025kubernetes, k8scontrollers, k8sclustercomponents}

\subsection{Extending Kubernetes with custom controllers}

The strength of Kubernetes here is the ability to define so-called \glspl{crd}.
These are themselves Kubernetes objects, which allow users of Kubernetes to specify custom resources in the Kubernetes \gls{api}.
These objects may also allow the management of resources other than container workloads, since the controllers attached to them simply need to read the desired state from Kubernetes and can then perform arbitrary actions, such as reaching out to external \glspl{api} or creating other Kubernetes resources.
\parencite{k8scustomresources, k8scontrollers}

The concept of writing controllers that manage other software or aggregate capabilities to orchestrate complex tasks is referred to as the \enquote{operator pattern}.
An \enquote{operator} here is a piece of software that is able to manage a complex task in the same way a human operator would.
Common examples of operator frameworks are Kubebuilder, Java Operator SDK and Kopf.
These frameworks aim to simplify the development of operators by reducing repetitive code, allowing developers to focus on business logic.
\parencite{k8soperators}

\subsubsection{Asynchronous reconciliation}
Controllers make use of a concept called asynchronous reconciliation.
It means that the kube-apiserver only verifies the syntactical correctness of the \texttt{CustomResource} created based on a \gls{crd}, but usually does not validate the actual logical consistency.
It is possible to implement such checks using admission webhooks, but validating logical consistency via admission webhooks is generally discouraged.
Validation of the actual logical consistency and the ability to move the observed state to the desired state should be determined at reconcile time.
This in turn means that while an object may be accepted by the kube-apiserver, it does not automatically mean that the desired state will be able to match the observed state of the cluster.
In promise theory terms, this means that the assessment of whether the promise of the controller has been kept cannot be made at admission time, but only over time, and only by observing the status reported by the controller on the object.
It also means that the client cannot assume that the promise has been fulfilled after the object has been accepted by the kube-apiserver, but that it may take a large amount of time, sometimes minutes or longer, to be fulfilled.
In Kubernetes, this concept is referred to as \enquote{eventually consistent}, the concept of simply waiting for the cluster to reach a consistent state where everything is running.
\parencite{k8soperators, k8scustomresources, k8scontrollers, k8sdynamicadmissioncontrol, k8sadmissioncontrol}

\subsubsection{Kubernetes object conditions}

Object conditions map to the assessments found in promise theory, representing the controller's view of whether it has fulfilled its promise.
Conditions, when following the convention, should use positive naming.
Typically, this amounts to a condition called \enquote{Ready} and additional conditions on more complex objects.
Every condition can then have additional properties, the primary one indicating whether the condition is active, with additional ones indicating a reason as to why the resource is in the specific condition, a human-readable message from the controller as well as the transition time.
\parencite{kubebuilder-book}

\subsubsection{Garbage collection}
Garbage collection is a concept making use of owner references to allow for the automatic deletion of objects when the managing resource is deleted.
For this to work, the controller sets an owner reference on the managed object pointing back to the parent object.
Owner references allow several features to work properly, including enabling the controller to be notified when a child object is modified.
The owner reference also tells the kube-apiserver to delete any child resources when the parent is deleted, allowing for easy cleanup of child resources without making use of finalizers.
\parencite{k8sgarbagecollection}
